\begin{tomtat}
	\textbf{ĐỀ 04}
\end{tomtat}
\begin{ex}
	Cho $a > 0$ và $a \ne 1$. Rút gọn biểu thức $P = \left(\frac{1}{\sqrt{a}+2} - \frac{1}{a+2\sqrt{a}}\right) : \frac{\sqrt{a}-1}{2\sqrt{a}+4}$.
	\loigiai{
		Với $a > 0$ và $a \ne 1$, ta có:
		\begin{align*}
			P &= \left(\frac{1}{\sqrt{a}+2} - \frac{1}{\sqrt{a}(\sqrt{a}+2)}\right) : \frac{\sqrt{a}-1}{2(\sqrt{a}+2)} \\
			&= \frac{\sqrt{a} - 1}{\sqrt{a}(\sqrt{a}+2)} : \frac{\sqrt{a}-1}{2(\sqrt{a}+2)} \\
			&= \frac{\sqrt{a} - 1}{\sqrt{a}(\sqrt{a}+2)} \cdot \frac{2(\sqrt{a}+2)}{\sqrt{a}-1}
		\end{align*}
		Vì $a \ne 1$ nên $\sqrt{a} \ne 1$, do đó $\sqrt{a}-1 \ne 0$. Ta có thể rút gọn $\sqrt{a}-1$.
		Vì $a > 0$ nên $\sqrt{a}+2 \ne 0$. Ta có thể rút gọn $\sqrt{a}+2$.
		\begin{align*}
			P &= \frac{1}{\sqrt{a}} \cdot \frac{2}{1} \\
			&= \frac{2}{\sqrt{a}}
		\end{align*}
		Vậy $P = \frac{2}{\sqrt{a}}$.
	}
\end{ex}

\begin{ex}
	Cho phương trình $x^2 - 4x + 1 = 0$. Không giải phương trình, hãy tính giá trị biểu thức $T = |x_1 - x_2|$.
	\loigiai{
		Phương trình đã cho là $x^2 - 4x + 1 = 0$.
		Đây là phương trình bậc hai có dạng $ax^2+bx+c=0$ với $a=1$, $b=-4$, $c=1$.
		Ta tính biệt thức $\Delta = b^2-4ac = (-4)^2 - 4(1)(1) = 16 - 4 = 12$.
		Vì $\Delta = 12 > 0$ nên phương trình có hai nghiệm phân biệt $x_1, x_2$.
		
		Theo hệ thức Vi-ét, ta có:
		$$ \begin{cases} 
			x_1+x_2 = \dfrac{-b}{a} = \dfrac{-(-4)}{1} = 4 \\ 
			x_1x_2 = \dfrac{c}{a} = \dfrac{1}{1} = 1 
		\end{cases} $$
		Ta cần tính giá trị biểu thức $T = |x_1 - x_2|$.
		Ta biết rằng $(x_1-x_2)^2 = (x_1+x_2)^2 - 4x_1x_2$.
		Do đó, $|x_1-x_2| = \sqrt{(x_1+x_2)^2 - 4x_1x_2}$ (vì $|x_1-x_2| \ge 0$).
		Thay giá trị $x_1+x_2$ và $x_1x_2$ đã tính ở trên vào biểu thức:
		\begin{align*}
			T &= \sqrt{(4)^2 - 4(1)} \\
			&= \sqrt{16 - 4} \\
			&= \sqrt{12} \\
			&= \sqrt{4 \cdot 3} \\
			&= 2\sqrt{3}
		\end{align*}
		Vậy giá trị của biểu thức $T$ là $2\sqrt{3}$.
	}
\end{ex}

\begin{ex}
	Giải hệ phương trình sau: $\begin{cases} x - 3y = 4 \\ 2x + y = 1 \end{cases}$.
	\loigiai{
		Ta có hệ phương trình:
		$$ \begin{cases} x - 3y = 4 & (1) \\ 2x + y = 1 & (2) \end{cases} $$
		Từ phương trình (1), ta biểu diễn $x$ theo $y$: $x = 4 + 3y$.
		Thay biểu thức của $x$ này vào phương trình (2):
		\begin{align*}
			2(4 + 3y) + y &= 1 \\
			8 + 6y + y &= 1 \\
			7y &= 1 - 8 \\
			7y &= -7 \\
			y &= -1
		\end{align*}
		Thay giá trị $y=-1$ vào biểu thức $x = 4 + 3y$:
		\begin{align*}
			x &= 4 + 3(-1) \\
			x &= 4 - 3 \\
			x &= 1
		\end{align*}
		Vậy nghiệm của hệ phương trình là $(x; y) = (1; -1)$.
	}
\end{ex}

\begin{ex}
	Trong trò chơi gieo thẻ từ thẻ số 1 đến thẻ số 8. Tính xác suất của biến cố "thẻ lấy ra được ghi số chia hết cho 4".
	\loigiai{
		Tập hợp tất cả các kết quả có thể xảy ra là:
		 $\{1, 2, 3, 4, 5, 6, 7, 8\}$. \\
		Tập hợp các kết quả thuận lợi cho biến cố "thẻ lấy ra được ghi số chia hết cho 4" là: $\{4, 8\}$. \\
		Xác suất của biến cố là: $P(A) = \frac{\text{Số kết quả thuận lợi}}{\text{Tổng số kết quả có thể}} = \frac{2}{8} = \frac{1}{4}$.
	}
\end{ex}

\begin{ex}
	Một mảnh vườn hình chữ nhật có chiều dài hơn chiều rộng 15 m. Nếu giảm chiều dài 2 m và tăng chiều rộng 3 m thì diện tích tăng thêm 44 m$^2$. Tính diện tích mảnh vườn đó.
	\loigiai{
		Gọi chiều dài của mảnh vườn hình chữ nhật ban đầu là $x$ (m) và chiều rộng là $y$ (m).
		Điều kiện: $x > 0, y > 0$ và $x > y$.
		
		Theo đề bài, chiều dài hơn chiều rộng 15 m, nên ta có phương trình thứ nhất:
		$$ x - y = 15 $$
		
		Khi giảm chiều dài 2 m,  tăng chiều rộng 3 m, thì diện tích tăng thêm 44 m$^2$, nên:
		\begin{align*}
			(xy + 3x - 2y - 6) - xy &= 44 \\
			3x - 2y - 6 &= 44 \\
			3x - 2y &= 50
		\end{align*}
		Ta có hệ phương trình:
		$ \begin{cases} x - y = 15 \\ 3x - 2y = 50 \end{cases} $
		\\
		Giải hệ phương trình ta được $x = 20$, $y = 5$ các giá trị này thõa mãn điều kiện ban đầu
		
		Chiều rộng ban đầu là $y = 5$ m.
		Chiều dài ban đầu là $x = 20$ m.
		Diện tích ban đầu của mảnh vườn là $S_1 = x \cdot y = 20 \cdot 5 = 100$ m$^2$.
		
		
		Vậy diện tích mảnh vườn đó là 100 m$^2$.
	}
\end{ex}




\begin{ex}
	Cho tam giác $ABC$ nhọn. Ba đường cao $AI, BK, CL$ cắt nhau tại $H$. Chứng minh:
	\begin{enumerate}
		\item Tứ giác $BIHL$ là tứ giác nội tiếp.
		%\item $\widehat{AKL} = \widehat{IKC}$ và $H$ là tâm đường tròn nội tiếp tam giác $IKL$.
	\end{enumerate}
	\loigiai{
		\begin{center}
			\begin{tikzpicture}[line join = round, line cap=round,>=stealth,font=\footnotesize,scale=1]
				% Def
				\def\r{3}
				
				% Coordinates
				\coordinate (O) at (0,0);
				\coordinate (A) at (110:\r);
				\coordinate (B) at (-160:\r);
				\coordinate (C) at (-20:\r);
				
				\coordinate (I) at ($(B)!(A)!(C)$); % I la hinh chieu cua A len BC
				\coordinate (L) at ($(A)!(C)!(B)$); % L la hinh chieu cua C len AB
				\coordinate (K) at ($(A)!(B)!(C)$); % K la hinh chieu cua B len AC
				
				% De tim H, can \usetikzlibrary{intersections} o preamble
				\path[name path=AI] (A) -- (I);
				\path[name path=CL] (C) -- (L);
				\coordinate[name intersections={of=AI and CL, by=H}]; % H la giao diem AI va CL
				
				% Draw
				\draw (A) -- (B) -- (C) -- cycle; % Ve tam giac ABC
				\draw (A) -- (I); % Ve duong cao AI
				\draw (C) -- (L); % Ve duong cao CL
				\draw (B) -- (K); % Ve duong cao BK
				\draw (O) circle (\r);
				% Ky hieu goc vuong bang lenh \pic
				% Can \usetikzlibrary{angles, quotes} o preamble
				\pic[draw, angle radius=2.5mm]{right angle = B--L--C}; % Goc BLC vuong
				\pic[draw, angle radius=2.5mm]{right angle = A--I--C}; % Goc AIC vuong (hoac AIB)
				\pic[draw, angle radius=2.5mm]{right angle = B--K--A}; % Goc BKA vuong (hoac CKA)
				
				% Fill 
				\fill (O) circle (1.5pt);
				\fill (A) circle (1.5pt);
				\fill (B) circle (1.5pt);
				\fill (C) circle (1.5pt);
				\fill (I) circle (1.5pt);
				\fill (L) circle (1.5pt);
				\fill (K) circle (1.5pt);
				\fill (H) circle (1.5pt);
				
				% Nodes	
				\node[above left] at (O) {$O$}; % Vi tri nhan co the can dieu chinh
				\node[above] at (A) {$A$};
				\node[below left] at (B) {$B$};
				\node[right] at (C) {$C$};
				\node[below right] at (I) {$I$}; % Vi tri nhan co the can dieu chinh
				\node[above left] at (L) {$L$}; % Vi tri nhan co the can dieu chinh
				\node[above right] at (K) {$K$}; % Vi tri nhan co the can dieu chinh
				\node[left] at (H) {$H$}; % Vi tri nhan co the can dieu chinh
				% Ve va danh dau doan thang bang nhau
			\end{tikzpicture}
			
		\end{center}
			\begin{enumerate}
			\item Tứ giác $BIHL$ là tứ giác nội tiếp.
			
			
			\item $\widehat{AKL} = \widehat{IKC}$ và $H$ là tâm đường tròn nội tiếp tam giác $IKL$.
			
			AI, CL là đường cao của tam giác ABC nên $AI \perp BC$ và $CL \perp AB$. Do đó
			$\widehat{AIB} = \widehat{BLC} = 90^\circ$ hay $\widehat{HIB} = \widehat{BLH} = 90^\circ$.
			
			hai điểm I, L cùng nằm trên đường tròn đường kính BH.
			
			Bốn điểm B, I, L, H cùng nằm trên đường tròn đường kính BH hay tứ giác BIHL nội tiếp đường tròn đường kính BH.
			
			Chứng minh tương tự câu 1, ta có tứ giác CIHK nội tiếp đường tròn đường kính CH.
			
			$\widehat{IKC} = \widehat{IHC}$ (hai góc nội tiếp cùng chắn cung IC)
			
			Chứng minh tương tự, ta có tứ giác AKHL nội tiếp đường tròn đường kính AH nên
			$\widehat{AKL} = \widehat{AHL}$ (hai góc nội tiếp cùng chắn cung AL).
			
			Lại có $\widehat{IHC} = \widehat{AHL}$ (đối đỉnh)
			
			Do đó $\widehat{AKL} = \widehat{IKC}$.
			
			+ Ta có $\widehat{AKL} + \widehat{LKB} = 90^\circ$ và $\widehat{IKC} + \widehat{IKB} = 90^\circ$
			
			Mà $\widehat{AKL} = \widehat{IKC}$ (câu 2) nên $\widehat{LKB} = \widehat{IKB}$ hay KB tức KH là tia phân giác của $\widehat{IKL}$.
			
			Chứng minh tương tự, ta có IH là tia phân giác của $\widehat{LIK}$.
			
			Xét tam giác IKL có KH, IH là hai đường phân giác của tam giác cắt nhau tại H nên H là tâm đường tròn nội tiếp tam giác IKL.
			
		\end{enumerate}
		}
\end{ex}